% ===== 原文 pp.55–58（PDF p.65–68） =====

% [原文 p.55]
$r,s,t=1,2,3$ 指 $W$ 的坐标）。对 3d 规范扇区，我们有
\begin{gather*}
    A_r = A_r^{[i]}, \\
    \varphi_r = \phi_r^{[i]}, \\
    \cA_r = (A_r+\ii\phi_r)^{[i]} \\
    \bar\cA_r = (A_r-\ii\phi_r)^{[i]} \\
    \lambda_r = \ii\left(\psi_r+\ii\tilde\psi_r\right)^{[i]}, \\
    \zeta_{rs} = -\ii\left(\chi^+-\ii\chi^-\right)_{rs}^{[i]}, \\
    \rho = \frac{1}{\sqrt2}\left(\psi_4-\ii\tilde\psi_4\right)^{[i]}, \\
    \tilde\rho = \half\left(\eta-\ii\tilde\eta\right)^{[i]} .
\end{gather*}
而物质扇区我们有
\begin{gather*}
    \sigma = \sqrt2\,\sigma^{[i]}, \\
    \tau = (A_4+\ii\phi_4)^{[i]}, \\
    \eta^{\bar\sigma} = \sqrt2\,\ii\left(\eta+\ii\tilde\eta\right)^{[i]}, \\
    \eta^{\bar\tau} = 2\ii\left(\psi_4+\ii\tilde\psi_4\right)^{[i]}, \\
    \chi_r^\sigma = \frac{1}{\sqrt2}\left(\psi_r-\ii\tilde\psi_r\right)^{[i]}, \\
    \chi_r^\tau = \left(\chi^+-\ii\chi^-\right)_{r4}^{[i]} .
\end{gather*}

% [Tutor Note: p.55 的两张表可按 (5.1) 逐条核对：由 $\delta_QA=\lambda$ 与 $\lambda_r=\ii(\psi_r+\ii\tilde\psi_r)^{[i]}$ 应反推出 $t=\ii$ 的 $\delta_QA_r=\ii(\psi_r+\ii\tilde\psi_r)$；由 $\delta_Q\varsigma=\rho$ 得 $\delta_Q\varphi_4=\frac{1}{\sqrt2}(\psi_4-\ii\tilde\psi_4)$；由 $\delta_Q\bar\sigma=\tilde\rho$ 得 $\delta_Q\bar\sigma=\half(\eta-\ii\tilde\eta)$。三处归一化若有一处写错（漏掉 $\sqrt2$ 或 $\ii$），这套“字典”会立刻矛盾，这是核对本表最快的办法。]

此外，引入一个辅助 0-形式 $P$ 以便使 BRST 变分在离壳时幂零是有用的；作用量中依赖于 $P$ 的项被选取来保证其运动方程为
\[
    P = \dd_A^\star\varphi - \frac{\ii}{2}\left([\bar\sigma,\sigma]+[\bar\tau,\tau]\right) .
\]

% [原文 p.56]
重新标记后的场的 BRST 变分为
\begin{gather*}
    \delta_QA = \lambda, \qquad \delta_Q\sigma = 0, \\
    \delta_Q\varphi = \ii\lambda, \qquad \delta_Q\tau = 0, \\
    \delta_Q\lambda = 0, \qquad \delta_Q\bar\sigma = \eta^{\bar\sigma} \\
    \delta_Q\zeta = -\ii\cF, \qquad \delta_Q\bar\tau = \eta^{\bar\tau} \\
    \delta_Q\rho = [\tau,\sigma], \qquad \delta_Q\eta^{\bar\sigma} = 0, \\
    \delta_Q\tilde\rho = \ii P, \qquad \delta_Q\eta^{\bar\tau} = 0, \\
    \delta_QP = 0, \qquad \delta_Q\chi^\sigma = \dd_{\cA}\sigma, \\
    \delta_Q\chi^\tau = \dd_{\cA}\tau .
\end{gather*}
其中 $\dd_{\cA}$ 与 $\cF$ 分别是关于复化联络 $\cA$ 的协变导数与曲率。用重新标记后的场写出的作用量为
\[
    S = \int_M\dd^4x\sqrt{g}\left(\cL_1+\cL_2+\cL_3\right) + S_{top}
\]
其中
\begin{align*}
    \cL_1 &= -\frac{1}{2e^2}\delta_Q\Tr\left\{\frac{\ii}{2}\zeta_{rs}\bar\cF^{rs}
        + \ii\tilde\rho\left(P - 2\dd_A^\star\varphi + \ii[\bar\sigma,\sigma] + \ii[\bar\tau,\tau]\right)\right\} \\
    \cL_2 &= -\frac{1}{2e^2}\delta_Q\Tr\left\{\chi_r^\tau\left(\dd_{\bar\cA}\bar\tau\right)^r
        + \chi_r^\sigma\left(\dd_{\bar\cA}\bar\sigma\right)^r\right\} \\
    \cL_3 &= -\frac{1}{2e^2}\delta_Q\Tr\left\{\rho\,[\bar\tau,\bar\sigma]\right\}, \\
    S_{top} &= \frac{\sqrt2}{e^2}\int_W\Tr\left\{\ii\chi^\tau\wedge \dd_{\cA}\chi^\sigma
        - \zeta\wedge\left(\dd_{\cA}\rho + [\chi^\sigma,\tau] - [\chi^\tau,\sigma]\right)\right\} .
\end{align*}

% [Tutor Note: p.56 的两个变分可用来检验 $\delta_Q^2=0$：$\delta_Q\chi^\tau=\dd_{\cA}\tau$ 且 $\delta_Q\tau=0$，故 $\delta_Q^2\chi^\tau=0$；又 $\cF=\dd_{\cA}\cA$ 且 $\delta_Q\cA=0$（因 $\delta_QA=\lambda$、$\delta_Q\varphi=\ii\lambda$ 给出 $\delta_Q(A+\ii\varphi)=0$），故 $\delta_Q^2\zeta=-\ii\delta_Q\cF=0$。这也顺带说明 $\cA=A+\ii\phi$ 是 BRST 不变的复化联络。]

把 $\tau$ 换成 $h_{\C}$ 之后，某些物质场切于（复化）群的非单位元点，用李代数取值的平均值来描述它们是方便的（正如上述 $t=1$ 约化那样）。也就是说，$\eta^{\bar\tau}$（取值于 $T_{\bar h_{\C}}$ 中）与 $\chi^\tau$（取值于 $T_{h_{\C}}$ 中）被换成平均后的场
\begin{align*}
    \hat\eta^{\bar\tau} &\equiv -\frac{1}{2\pi}\int_0^{2\pi}\dd x^4\, \bar H_{\C}^{-1}\eta^{\bar\tau}\bar H_{\C} \\
    \hat\chi^\tau &\equiv -\frac{1}{2\pi}\int_0^{2\pi}\dd x^4\, H_{\C}^{-1}\chi^\tau H_{\C} .
\end{align*}

% [原文 p.57]
用新变量写出的 BRST 变分如下
\[
    \delta_Qh_{\C} = 0, \qquad
    \delta_Q\bar h_{\C} = 2\pi\bar h_{\C}\hat\eta^{\bar\tau}, \qquad
    \delta_Q\hat\chi^\tau = (2\pi)^{-1}h_{\C}^{-1}\dd_{\cA}h_{\C}, \qquad
    \delta_Q\hat\eta^{\bar\tau} = 0 .
\]
约化步骤与我们对 Vafa-Witten、$t=0$、$t=1$ 理论所做过的类似：我们用一个 BRST-正合的量形变作用量，以便用新变量写出它，然后通过令场不依赖 $x^4$ 坐标而约化到 3d。所得的 3d 作用量为
\[
    \hat S = \int_W\dd^3x\sqrt{g_W}\left(\hat\cL_1+\hat\cL_2+\hat\cL_3\right) + S_{top}
\]
其中
\begin{align*}
    \hat\cL_1 &= -\frac{1}{2e^2}\delta_Q\Tr\left\{\frac{\ii}{2}\zeta_{rs}\bar\cF^{rs}
        + \ii\tilde\rho\left(P - 2\dd_A^\star\varphi + \ii[\bar\sigma,\sigma] + \ii[\bar\tau,\tau] + \cdots\right)\right\} \\
    \hat\cL_2 &= -\frac{1}{2e^2}\delta_Q\Tr\left\{-(2\pi)^{-1}\hat\chi_r^\tau\,\bar h_{\C}^{-1}
        \left(\dd_{\bar\cA}\bar h_{\C}\right)^r + \chi_r^\sigma\left(\dd_{\bar\cA}\bar\sigma\right)^r\right\} \\
    \hat\cL_3 &= -\frac{1}{2e^2}\delta_Q\Tr\left\{-(2\pi)^{-1}\rho\,\bar h_{\C}^{-1}[\bar h_{\C},\bar\sigma]\right\},
\end{align*}
其中 $\hat\cL_1$ 中的 $\cdots$ 表示我们未显式写出的高阶项。靶空间上的 $G$-不变辛形式正比于
\[
    \Omega = \Tr \dd\sigma\, h_{\C}^{-1}\dd h_{\C} .
\]
靶空间上的 $G$-不变 Kähler 形式对应于 [63] 中描述的超 Kähler 度规，且不容许简单的闭式表达式；在单位元的邻域内它正比于
\[
    J = \frac{\ii}{2}\Tr\left(\dd\sigma\dd\bar\sigma + \dd\tau\dd\bar\tau\right) .
\]
$G$-作用关于辛形式 $J$、$\Omega$、$\bar\Omega$ 的矩映射 $\mu_3,\mu_+,\mu_-$ 正比于坐标的如下函数：
\begin{align*}
    \mu_3 &= -\frac{\ii}{2}\left([\bar\sigma,\sigma]+[\bar\tau,\tau]+\cdots\right) \\
    \mu_+ &= -\ii h_{\C}^{-1}[h_{\C},\sigma] \\
    \mu_- &= -\ii \bar h_{\C}^{-1}[\bar h_{\C},\bar\sigma] .
\end{align*}

% [Tutor Note: p.57 的矩映射可与 p.55 的辅助场运动方程对账：$P=\dd_A^\star\varphi-\frac{\ii}{2}([\bar\sigma,\sigma]+[\bar\tau,\tau])$ 与 $\mu_3=-\frac{\ii}{2}([\bar\sigma,\sigma]+[\bar\tau,\tau]+\cdots)$ 恰给出 $P=\dd_A^\star\varphi+\mu_3$（忽略 $\cdots$ 的高阶项）。这是检验 $\mu_3$ 符号的最快办法。]
% [Tutor Note: p.57 的 $\Omega=\Tr\dd\sigma h_{\C}^{-1}\dd h_{\C}$ 是 $T^*G_{\C}$ 上的全纯辛形式：纤维坐标 $\sigma$ 与底坐标的 Maurer–Cartan 形式 $h_{\C}^{-1}\dd h_{\C}$ 配对，故它是 $T^*G_{\C}$ 的标准辛形式（与 $J$ 构成超 Kähler 三元组），可与 $X=T^*G_{\C}$ 的维数 $2\dim G$ 一起核对。]

% [原文 p.58]
我们已经证明，作用量与变分的 BRST-正合部分约化为 B.5 节所写的、靶空间为 $T^*[2]G_{\C}$ 的带规范场 Rozansky-Witten 模型的相应部分。我们猜测拓扑项 $S_{top}$ 也适当地约化为该模型的拓扑项（尽管我们在此不予证明）。
